\section{Construction}\label{sec:txid-construction}

The protocol proceeds as follows. The ledger operator processes user
transactions. Each user is provisioned with a long-term signing key and a
long-term tracing key, cryptographically bound together via a credential issued
by the registration authority. For each epoch, the user derives an epoch tracing
key from their tracing key and the epoch identifier. Transactions submitted
during that epoch are tagged using PRF $\FT$, evaluated on the epoch tracing key
and an index in the range $[0,\Max)$, where $\Max$ bounds the number of
transactions a user can submit per epoch. Each transaction carries a proof that
its tag was correctly generated.

For tracing, the user discloses their epoch tracing key to the relevant auditor,
who locally computes every valid transaction tag for that user and epoch. This
requires the ledger to be indexed by tag, so that the auditor pays one lookup
per candidate tag rather than a scan. We assume throughout that the ledger
operator maintains such an index, which is why tracing costs $O(\Max)$ rather
than $O(|L|)$ for a ledger of $|L|$ transactions.

When the system moves to a new epoch, users derive a fresh epoch tracing key,
and tracing capabilities granted for prior epochs expire automatically.

\paragraph{The index is unordered.} The index $\gamma \in [0,\Max)$ is the
counterpart of the monotonically increasing counter used by prior tag-based
schemes~\cite{CCS:KiaKohSar22,CANS:SarKiaKoh24}, and the difference between the
two is the point of the construction. Those schemes need the counter to increase
by exactly one per transaction, with no gaps. Their auditor walks the counter
upwards from zero and stops at the first miss, so a gap would hide everything
after it. That forces the user to serialise: the $(i+1)$-th transaction cannot
be built until the $i$-th is fixed. Here the auditor
evaluates $\FT$ at every $\gamma \in [0,\Max)$ regardless, so gaps and
out-of-order use are free. The only constraint is that a user must not reuse an
index within an epoch, and they can enforce that themselves. A user may
partition $[0,\Max)$ among concurrent workers in advance, or draw indices from a
local pool. Either way the resulting transactions can be submitted in any order,
and in parallel.

Sampling $\gamma$ uniformly instead is tempting but worse. Collisions start
appearing after roughly $\sqrt{\Max}$ transactions. A collision costs a rejected
transaction rather than a security failure, so nothing breaks, but explicit
allocation avoids the problem for free.

Once a user has consumed all $\Max$ indices they wait for the next epoch. So
$\Max$ is a per-user, per-epoch rate limit, and choosing it trades that limit
against the cost of tracing and the size of a first grant.

\paragraph{Building blocks.} The protocol is built from five cryptographic
primitives, each used abstractly. The registration authority certifies users'
keys with a digital signature scheme $\DS$.  Users bind their tracing keys
across several proofs using a commitment scheme $\Com$. Two pseudorandom
functions, $\FE$ for epoch keys and $\FT$ for tags, derive epoch tracing keys
and transaction tags from users' long-term tracing keys. These keys are also
bound using one-way functions $\G$ and $\G'$, where $\G'$ is applied to the
tracing key as a whole, which the instantiation takes to be a tuple. The image
$\PTK = \G'(\utk)$ is the user's public tracing key, and doubles as a public
parameter of $\FT$. The user
proves that these values were derived correctly, and, where the proof must also
authenticate a message, that the message originated from the same user, using
zero-knowledge proofs of knowledge.

Two requirements on these primitives are not generic, and the instantiation
of~\Cref{sec:txid-instantiation} depends on both. First, $\DS$ and $\Com$ must
be compatible, meaning that $\DS$ supports signatures on committed values. Two
things have to work. The registration authority must be able to certify
$\G(\sk)$ and $\G'(\utk)$ without learning $\sk$ or $\utk$, and a user must
later be able to prove in zero knowledge that a commitment $\com$ opens to the
same $\utk$ a credential certifies. Pick $\DS$ and $\Com$ independently and
$\rel_6$ below needs a generic proof over the signature verification equation,
which is the one expense the instantiation is built to avoid. Second, $\FT$ is
parameterised
by the public tracing key, as described next.

\paragraph{The public tracing key is a parameter of $\FT$.} The tag function
takes three arguments, $\Tag_\gamma = \FT(\PTK, \ek, \gamma)$, where $\PTK =
\G'(\utk)$ is the public tracing key the registration authority certifies at
registration. $\PTK$ enters as a public parameter and not as key material: the
key of $\FT$ is the epoch tracing key $\ek$, and $\PTK$ is fixed once, for the
lifetime of the user. Correspondingly, we require pseudorandomness of both
families to hold against a distinguisher that holds $\PTK$, in the sense
of~\Cref{sec:prelims-prfs}, so that nothing $\PTK$ carries need be kept secret.

The distinction matters because whatever the instantiation needs in order to
evaluate a tag, beyond the epoch key itself, must sit somewhere that does not
vary with the epoch. Placing such material in the output of $\FE$ would make
$\FE$ repeat a component of its output in every epoch, and no function that does
that is pseudorandom, whatever the hardness assumption behind it. Placing it in
$\PTK$ costs nothing, since $\PTK$ is certified at registration and already
travels with the grant message, and it leaves $\FE$ a pseudorandom function in
the ordinary sense. \Cref{sec:txid-instantiation} is exactly this case: $\PTK$
carries a vector of $\Max$ group elements alongside the image of $\utk$, while
$\ek$ is a single group element.

\paragraph{Setup.} During setup, the registration authority generates a key pair
for a digital signature scheme. Their public key is communicated to the ledger
operator. The ledger operator generates the public parameters for a commitment
scheme and for four proof systems, used respectively for registration, for epoch
tracing key derivation, for transaction authorisation, and for tag correctness.

The registration authority generates $(\rask, \rapk) \leftarrow
\DSgen(\secparam)$ and communicates $\rapk$ to the ledger operator.  The ledger
operator calls $\ck \leftarrow \ComGen(\secparam, 1)$ and produces four sets of
public parameters, one for each relation.

The first relation, $\rel_4$, lets a user prove knowledge of their long-term
signing and tracing keys. Its instance, witness, and relation are as follows:
%
\begin{equation}
%
  \inst_4 = (\G, \G', \PSK, \PTK), \qquad \witness_4 = (\sk, \utk),
%
\end{equation}
%
\begin{equation}\label{eq:rel4} \begin{gathered}
%
  \rel_4 = \{\ \PSK = \G(\sk) \ \wedge \ \PTK = \G'(\utk)\ \}.
%
\end{gathered} \end{equation}
%
The second relation, $\rel_5$, captures correctness of epoch tracing key
generation, with instance, witness, and relation given as follows:
%
\begin{equation}
%
  \inst_5 = (\FE, \G', \PTK, \epoch, \ek), \qquad \witness_5 = \utk,
%
\end{equation}
%
\begin{equation}\label{eq:rel5} \begin{gathered}
%
  \rel_5 = \{\ \PTK = \G'(\utk) \ \wedge \ \ek = \FE(\utk, \epoch)\ \}.
%
\end{gathered} \end{equation}
%
Since $\PTK$ appears in $\inst_5$, the same proof establishes that the epoch key
was derived from the tracing key whose image the credential certifies, and that
$\PTK$ itself is well formed against that tracing key. In the instantiation the
second of these is proved once, at an auditor's first grant, and the first once
per epoch.
%
The third relation, $\rel_6$, does two jobs at once. It lets a user prove
possession of a valid registration credential, and prove that the tracing key
certified by that credential is the one inside a commitment. Its instance,
witness, and relation are as follows:
%
\begin{equation}
%
  \inst_6 = (\G, \G', \rapk, \ck, \com), \qquad \witness_6 = (\signature, \sk,
\utk, \openinfo),
%
\end{equation}
%
\begin{equation}\label{eq:rel6} \begin{gathered}
%
  \rel_6 = \{\ \DSverify(\rapk, \G(\sk) \| \G'(\utk), \signature) = 1 \ \wedge \
\ComOpen(\ck, \com, \openinfo) = \utk\ \}.
%
\end{gathered} \end{equation}
%
The fourth relation, $\rel_7$, lets a user prove that a given epoch tracing key
and transaction tag were correctly derived from the tracing key contained in the
same commitment, with instance, witness, and relation defined as follows:
%
\begin{equation}
%
  \inst_7 = (\FE, \FT, \G', \ck, \com, \epoch, \Max, \Tag_\gamma), \qquad
\witness_7 = (\utk, \openinfo, \ek, \gamma),
%
\end{equation}
%
\begin{equation}\label{eq:rel7} \begin{gathered}
%
  \rel_7 = \{\ \ComOpen(\ck, \com, \openinfo) = \utk \ \wedge \ \ek =
\FE(\utk, \epoch) \\ \wedge \ \Tag_\gamma = \FT(\G'(\utk), \ek, \gamma) \ \wedge
\ 0 \leq \gamma < \Max\ \}.
%
\end{gathered} \end{equation}
%
The ledger operator then publishes the following:
%
\begin{equation*} \ledger[\crs] \leftarrow (\G, \G', \FE, \FT, \rapk, \ck,
\crs_4, \crs_5, \crs_6, \crs_7, \Max).  \end{equation*}

\paragraph{Why four relations.} The number of relations reflects the roles the
ledger imposes, rather than an accumulation of unrelated design choices. Unlike
a credential system where a single verifier checks one presentation
interactively, transactions here are posted to a public, append-only ledger
whose transaction validators must accept or reject each one locally and
non-interactively, without contacting the sender, the registration authority, or
any auditor. $\rel_4$ certifies a user's long-term keys once, at registration.
$\rel_5$ re-derives a fresh, epoch-bound tracing capability that a user can
later hand to an auditor. $\rel_6$ and $\rel_7$ are the two checks the
transaction validators themselves run on every submitted transaction: that it
originates from a registered user, and that its tag is correctly derived.
Splitting the protocol this way is what lets validators, auditors, and users
each act without waiting on one another.

\paragraph{Registration.} During registration, the user proves knowledge of
their long-term signing and tracing keys, and the registration authority
certifies the corresponding public values, binding them together for future use
without ever learning the keys themselves.

The registration authority first checks that the user has not registered before,
then sends a fresh nonce. The user generates a signing key pair and a tracing
key, derives their corresponding public values $\PSK$ and $\PTK$, and responds
with a zero-knowledge proof of knowledge of both secrets, bound to the nonce so
the proof cannot be replayed, exactly as captured by $\rel_4$.

If the proof verifies, the registration authority certifies the user's public
values with $\DSsign$, records the user as registered, and returns the resulting
signature. Together with the public values, this signature becomes the user's
registration credential.

\paragraph{Authorisation.} Before the start of epoch $\epoch$, an auditor
requests tracing authorisation from a user, identifying themselves and supplying
a fresh nonce. In response, the user derives their epoch tracing key from their
long-term tracing key and the epoch identifier, and proves that this derivation
is correct, exactly as captured by $\rel_5$. Since this proof does not depend on
the auditor, it can be generated once and reused for every authorisation granted
during the epoch.

To authorise this particular auditor, the user binds the epoch key proof and the
auditor's nonce together with a second zero-knowledge proof of knowledge of
their long-term keys, again relative to $\rel_4$, preventing the authorisation
from being replayed elsewhere. The user then sends the auditor their
registration credential, the epoch tracing key, and both proofs. The auditor
accepts the key once the credential verifies with $\DSverify$ and both proofs
check out.

The credential is sent to the auditor as issued, rather than randomised as it is
at submission time. This is deliberate, and it costs nothing. A grant is
addressed to a named auditor for a named user, so there is no identity left to
hide from them. And an auditor who replays the credential elsewhere cannot
produce the accompanying $\rel_4$ proof, which is bound to a fresh nonce and
needs knowledge of $\sk$.

\paragraph{Transaction submission.} During transaction submission, the user
proves that the transaction tag was correctly generated from their tracing key,
while also authenticating the transaction itself.

The user prepares a confidential payload and picks a fresh, unused index within
the epoch's allowed range. They derive the transaction tag from their epoch
tracing key and that index, and commit afresh to their long-term tracing key,
since the checks that follow relate the tag back to a committed key rather than
to the key itself.

Two proofs accompany the transaction. The first shows that the tag and the fresh
commitment are consistent with each other and with a single underlying tracing
key, and that the index lies in $[0,\Max)$, exactly as captured by $\rel_7$. The
second binds the whole transaction, the payload, the tag, the commitment, and
the first proof, to a valid registration credential for that same committed key,
as captured by $\rel_6$. Since this proof must authenticate a message rather
than merely certify correctness, it is generated by running $\NIZKprove$ on the
rest of the transaction as the message being signed. The user submits the
resulting tuple to the ledger.

The transaction validators run $\algo{Verify}$, rejecting the transaction
outright if its tag has already accompanied a previously accepted transaction,
since a repeated tag would let two transactions share an index and become
linkable. Otherwise, they accept the transaction once both proofs verify with
$\NIZKverify$, one against $\rel_7$ and the other against $\rel_6$, and record
the payload against the tag.

\paragraph{The range clause is load-bearing.} The clause $0 \leq \gamma < \Max$
in $\rel_7$ is not a formality. An auditor holding $\ek$ evaluates $\FT$ only at
indices in $[0,\Max)$, so a user who could tag a transaction with $\gamma =
\Max$ would obtain a transaction that no permission can ever trace, defeating
non-repudiation. Together with the validators' tag-uniqueness check, which is
what stops a user reusing an index, the range clause is what makes
$|\txns[(\user,\epoch)]| \leq \Max$ hold in the real world and match step 3 of
$\IDsubmit$ in $\idfu{\algo{trace}}$. Both checks must therefore be enforced by
the validators on every transaction.

\paragraph{Tracing.} Given the user's public tracing key $\PTK$ and epoch key
$\ek$ for epoch $\epoch$, the auditor traces the transactions the user submitted
during that epoch by computing $\Tag_\gamma = \FT(\PTK, \ek, \gamma)$ for each
index $\gamma \in [0,\Max)$ and retrieving entry $\ledger[\Tag_\gamma]$ where
present. Both values reach the auditor in the grant message of the epoch, so
tracing requires no state carried over from any earlier epoch.

\paragraph{Increment Epoch.} The system participants move to epoch $\epoch + 1$.

\begin{theorem}\label{thm:txid-bb} Assuming that $\DS$ satisfies existential
unforgeability, $\FE$ and $\FT$ are pseudorandom against a distinguisher holding
$\PTK$, $\G$ and $\G'$ are one way, and the zero-knowledge proof systems for
$\rel_4$ through $\rel_7$
are zero knowledge and simulation extractable, the protocol described above
securely realises $\idfu{\algo{trace}}$.  \end{theorem}
